\scanpage{73}
% archive sheet number: 72

[Ici, on montre que \inserted{la condition a) implique que} pour tt $X$ fixé, l'expression de a), $G(f,g)=R(f,f',\ldots,f^{(n)};g,\unsure{g|g'},\ldots,g^{(n)})$ doit avoir des coefficients $\in\mathbb C$, indépendants de $X$, donc

\[
R\in\mathbb C(F_0,F_1,\ldots,F_n;G_0,\ldots,G_n)
\]

et en fait, $R(\quad)=\underbrace{P(F_0,\ldots;G_0,\ldots)}_{\in\mathbb C[F_0,\ldots,G_0\ldots]}/F_1^p G_1^q$ \hfill ]

On peut prendre

\[
G(f,g)=\biggl(\left[-2\,\frac{g'''}{g'}+3\left(\frac{g''}{g'}\right)^2\right]-
 \left[-2\,\frac{f'''}{f'}+3\left(\frac{f''}{f'}\right)^2\right]dt^2
\]

De plus, on trouve un torseur canonique $\mathcal T_X$ \deleted{\unsure{sous|sur}} $\underline\omega^{\otimes2}_{X/\mathbb C}$, \unsure{dépend|dépent|dépendant} fonct\unsure{ionellt|ionellement|oriellt|oriellement} de $X$ (pour morphismes étales), et

\[
\deleted{\unsure{c_0|G_0}}\ c^X:\mathcal P_X\longrightarrow\mathcal T_X
\]

tel que

\[
\left\{\begin{aligned}
G(f,g)&=c^X(g)-c^X(f)\\
c^X&\text{ se factorise en } \!{}_H\backslash\mathcal P_X\longrightarrow\mathcal T_X\\
&\text{Ce dernier morphisme est un iso}
\end{aligned}\right.
\]

\uline{III}) On prend maintenant $X$ \inserted{schéma} variable lisse (ou seulement diff. lisse à l'ordre 3) sur des schémas variables, et on cherche pour tt tel $X$

\[
\unsure{\mathcal E_X|\mathscr E_X}\times\mathcal E_X\xrightarrow{G}\omega^2_{X/S}
\]

de façon compatible avec chgt de base $S'\to S$, avec morphismes étales

\diagram{\begin{tikzcd}[ampersand replacement=\&,column sep=small,row sep=small]
X'\arrow[rr]\arrow[dr,dash] \& \& X\arrow[dl,dash]\\
\& S \&
\end{tikzcd}}

, et satisfaisant les conditions qui correspondent à \inserted{$\Rightarrow$ b), c),} (\uline{plus} les condition d) \deleted{\unsure{e)}} \inserted{si \deleted{\unsure{on|si}} $S$ corps de car. 0, et e) dans} \deleted{\unsure{b), c) du cas précédent de l'} analytique complexe}).

% slanted
\margin{la forme suivante : $G(f,g)$ n'est \unsure{identiq$^{\mathrm{mt}}$|ident.} nul en aucune carct. $p>0$}

Mais quand on parle d'``op. différentiel'', il faut maintenant remplacer un polynôme $R$ en $f,f',f'',\ldots,g,g',g'',g'''$ par un polynôme en $f,f'=D_t^{(1)}f,D_t^{(2)}(f)\ldots D_t^{(n)}f,\ g,g'=D_t^{(1)}g,D_t^{(2)}g,\ldots D_t^{(n)}g$. \deleted{On voit comme} A priori, les coeff. sont dans $\Gamma(U,\underline{\mathcal O}_U)$. (si $t$, \deleted{\unsure{est|de}} l'uniformisante, est définie sur $U$) mais si on \deleted{\unsure{f|p}} veut faire varier $X$, on voit que les coefficients doivent être \uline{dans $\mathbb Z$} (ou dans

$\Gamma(S_0,\underline{\mathcal O}_{S_0})$ si on a fixé un schéma de base commun $S_0$ pour les $S,X\ldots$). On trouve encore l'existence d'un tel, et l'unicité à un facteur $\varepsilon\in\mathbb Z^*$ près, i.e. $\varepsilon=\pm1$ : \uline{unicité au signe près} ! Ici, on ne peut prendre l'expression $G'(f,g)$ précédente, car elle est identiq$^{\mathrm{mt}}$ nulle en car. 2 et 3 ! Il faut prendre $\frac1{12}G(f,g)$, savoir

\[
G_0(f,g)=\left[\left(-\frac{D_t^{(3)}g}{D_t^{(1)}g}+\left(\frac{D_t^{(2)}g}{D_t^{(1)}g}\right)^2\right)
-\left(-\frac{D_t^{(3)}f}{D_t^{(1)}f}+\left(\frac{D_t^{(2)}f}{D_t^{(1)}f}\right)^2\right)\right]dt^2
\]

Pour caractériser ce choix (parmi les deux possibles au signe près) on va calculer p.ex

\[
G_0(t,t^2)=\deleted{\unsure{\frac1{2t}dt^2|\frac1{2}dt^2|\frac1{2t}}}\ \frac1{2t}dt^2\qquad
\text{(au lieu de }-\tfrac1{2t}\,dt^2\text{)}
\]

[Il ne faut pas s'effrayer du dénominateur, pour le cas où 2 n'est pas inv. sur $X$ : alors $t^2$ n'est jamais une uniformisante !] ou bien soit quelque chose qui ait un sens aussi en car 2

\[
G_0(t,t^3)=\underbrace{\left(-\frac1{3t^2}+\left(\frac{3t}{3t^2}\right)^2\right)}_{\frac1{t^2}(-\frac13+1)}dt^2
=\deleted{\unsure{\frac1{3t^2}|\frac1{3t}}}\ \frac23\,\frac1{t^2}dt^2
\]

[mais c'est inutile, car en car 2 les \deleted{\unsure{définit$^{\mathrm{ns}}$|définitions|différents}} \inserted{deux choix} possibles sur $\mathbb Z$ coïncident !].

\uline{Rem.} 1) On peut se demander quel est le rôle joué dans cet énoncé par les biformes, et s'il n'y aurait pas moyen d'éliminer ce rôle. On est alors conduit, dans I), à se donner un faisceau analytique coh. $\mathcal F$ sur $X$, et à supposer $G'\deleted{\unsure{(f,g)|f,g}}:\mathcal E_X\times\mathcal E_X\longrightarrow\mathcal F$. On trouve alors qu'un tel $G'$ s'écrit de façon unique (où $G$ est comme dans l'énoncé initial, p.ex. celui de II) comme $G'=uG$, où $u:\underline\omega_X^{\otimes2}\longrightarrow\mathcal F$, et $u$ est même un \uline{iso} ! Donc $\underline\omega_X^{\otimes2}$ joue ici un rôle \uline{universel}. Dans II) on pourrait se donner un $\mathcal F_X$ dépd$^{\mathrm{t}}$
